\chapter{Financement~: repo, prêt de titres et prime brokerage}\label{ch:m1:financing}

Un fonds détient des actions pour cinq milliards de dollars et dispose
d'un milliard de fonds propres. Les quatre autres sont empruntés, et le prêt est renouvelé chaque nuit par
des prêteurs qui peuvent refuser de le renouveler. Un jour ordinaire, personne
n'y pense. Le jour où les actions chutent de quatre pour cent, un cinquième des
fonds propres a disparu, ses prêteurs exigent davantage de collatéral d'ici l'après-midi, et la seule
façon de le trouver est de vendre pour 800~millions de dollars des mêmes actions sur un marché
en baisse. Presque toute position de la finance professionnelle est financée par
l'argent d'autrui. Ce chapitre traite de l'origine de cet argent,
de son coût, des trois contrats par lesquels il est prêté --- et de ce qui
se produit lorsqu'il s'en va.

\section{Effet de levier}

\begin{definition}[Effet de levier]\label{def:m1:financing:leverage}
Le \emph{effet de levier}\index{effet de levier} d'un portefeuille est le rapport $L = A/E$
entre les actifs qu'il détient et les fonds propres (capitaux propres) qui les
portent. Le reste, $A - E$, est financé.
\end{definition}

\begin{proposition}[Arithmétique du levier]\label{prop:m1:financing:roe}
Soit $r_A$ le rendement des actifs sur une période et $r_f$ le coût du financement. Alors
\begin{enumerate}
\item le rendement des fonds propres est $r_E = L\,r_A - (L-1)\,r_f$~;
\item les fonds propres sont anéantis lorsque les actifs baissent de $1/L$~;
\item pour rétablir le ratio de levier après une baisse des actifs d'une fraction $x$
(avec $x < 1/L$), le portefeuille doit vendre des actifs valant $(L-1)\,x\,A$.
\end{enumerate}
\end{proposition}
\begin{proof}
(1) Les fonds propres rapportent $A r_A - (A-E) r_f$~; on divise par $E$. (2) La perte $xA$
égale $E$ lorsque $x = E/A$. (3) Après la baisse, les actifs valent $A(1-x)$ et les
fonds propres valent $E - xA$. Les actifs compatibles avec le levier $L$ valent $L(E - xA) = A - LxA$.
L'excédent à vendre est $A(1-x) - A + LxA = (L-1)xA$.
\end{proof}

\begin{example}[Le fonds du paragraphe d'introduction]\label{ex:m1:financing:fund}
$A = \$5$ milliards, $E = \$1$ milliard, $L = 5$. Avec $r_f = 4\%$, une hausse de $+10\%$
sur les actifs donne $+34\%$ sur les fonds propres~; une année à $-10\%$ donne $-66\%$. Une
baisse de 20~\% suffit à mettre fin au fonds. Après une baisse de 4~\%, le fonds a perdu 200~millions de dollars,
soit un cinquième de ses fonds propres, et doit vendre $4 \times 0.04 \times 5 =
\$0.8$ milliard pour retrouver un levier de cinq~: \emph{quatre dollars de
vente pour chaque dollar perdu}.
\end{example}

\begin{omfigure}
\begin{tikzpicture}
\begin{axis}[width=10.5cm,height=5cm,
  xlabel={rendement des actifs (\%)},ylabel={rendement des fonds propres (\%)},
  tick label style={font=\small},label style={font=\small},
  xmin=-20,xmax=20,ymin=-110,ymax=140,grid=major,grid style={gray!25},
  legend columns=3,legend style={font=\footnotesize,at={(0.5,-0.32)},anchor=north,draw=none,fill=none,/tikz/every even column/.append style={column sep=8pt}},
  table/col sep=comma]
\addplot[omMeth,very thick] table[x=asset_ret_pct,y=l1]{figdata/markets-1/06-financing/roe.csv};
\addplot[omDef,very thick] table[x=asset_ret_pct,y=l3]{figdata/markets-1/06-financing/roe.csv};
\addplot[omThm,very thick] table[x=asset_ret_pct,y=l8]{figdata/markets-1/06-financing/roe.csv};
\addplot[black,dashed] coordinates {(-20,-100) (20,-100)};
\legend{$L=1$,$L=3$,$L=8$}
\end{axis}
\end{tikzpicture}

\omcaption{Rendement des fonds propres en fonction du rendement des actifs, avec un
financement au taux de 4~\%. La ligne en pointillés est la perte totale~: pour $L = 8$, elle est atteinte lorsque les
actifs baissent de 9~\% --- moins que $1/L = 12.5\%$, car une année de coût de financement
est aussi due. Données~: la \cref{prop:m1:financing:roe}, calculées par le
script du chapitre.}
\end{omfigure}

Qui prête le $A - E$~? Trois contrats font presque tout le travail~: le repo,
le \omterm{def:m1:financing:seclending}{prêt de titres} et le swap. Tous les trois sont \emph{garantis}~: le prêteur
détient un collatéral d'une valeur supérieure au prêt, et cet excédent est la grandeur
à surveiller.

\section{Repo}

\begin{definition}[Pension livrée et décote]\label{def:m1:financing:repo}
Dans une \emph{pension livrée}\index{pension livrée} (repo), une partie
vend un titre contre des espèces et s'engage dans le même temps à le racheter
à une date ultérieure fixée et à un prix fixe plus élevé. Économiquement, c'est un prêt
d'argent garanti par le titre~; l'écart de prix représente les intérêts, cotés comme le
\emph{taux de repo}. La \emph{décote}\index{décote} $h$ est la fraction dont
la somme prêtée est inférieure à la valeur de marché du titre~: un titre
valant 100 permet de lever $100(1-h)$.
\end{definition}

\begin{omfigure}
\begin{tikzpicture}[font=\small,
  box/.style={draw=omDef,rounded corners=2pt,minimum width=3cm,minimum height=1cm,align=center,fill=omDef!4},
  flow/.style={-{Stealth[length=2mm]},thick}]
\node[box] (b1) at (0,1.5) {Emprunteur de liquidités\\(propriétaire du titre)};
\node[box] (l1) at (9,1.5) {Prêteur de liquidités};
\draw[flow,omDef] ([yshift=5pt]b1.east) -- node[above,font=\footnotesize]{obligation valant 100} ([yshift=5pt]l1.west);
\draw[flow,omThm] ([yshift=-5pt]l1.west) -- node[below,font=\footnotesize]{liquidités 98 \quad (décote 2~\%)} ([yshift=-5pt]b1.east);
\node[font=\footnotesize\bfseries,anchor=east] at (-2.0,1.5) {Aujourd'hui};
\node[box] (b2) at (0,-1.1) {Emprunteur de liquidités};
\node[box] (l2) at (9,-1.1) {Prêteur de liquidités};
\draw[flow,omThm] ([yshift=5pt]b2.east) -- node[above,font=\footnotesize]{liquidités $98\,(1 + r\,\tfrac{n}{360})$} ([yshift=5pt]l2.west);
\draw[flow,omDef] ([yshift=-5pt]l2.west) -- node[below,font=\footnotesize]{la même obligation} ([yshift=-5pt]b2.east);
\node[font=\footnotesize\bfseries,anchor=east] at (-2.0,-1.1) {Dans $n$ jours};
\end{tikzpicture}

\omcaption{Les deux temps d'un repo. Si l'emprunteur ne rachète pas, le
prêteur détient une obligation valant 100 face à une créance de 98~: la \omterm{def:m1:financing:repo}{décote} est
sa protection contre une baisse du cours de l'obligation pendant qu'il la vend.}
\end{omfigure}

\begin{proposition}[Les décotes bornent le levier]\label{prop:m1:financing:haircut}
Un portefeuille financé entièrement par repo à \omterm{def:m1:financing:repo}{décote} $h$ peut atteindre un levier
d'au plus $1/h$. Si la \omterm{def:m1:financing:repo}{décote} passe de $h$ à $h'$ sans modification des
prix, un portefeuille au levier maximal doit vendre une fraction $1 - h/h'$ de
ses actifs.
\end{proposition}
\begin{proof}
Chaque unité d'actifs permet de lever $1-h$ en espèces, donc les fonds propres nécessaires valent $hA$ et
$L = A/E \le 1/h$. À fonds propres inchangés, soit $E = hA$, le nouveau maximum est
$E/h' = (h/h')A$.
\end{proof}

\begin{example}[Deux pour cent à quatre pour cent]\label{ex:m1:financing:doubling}
Des obligations d'État financées avec une \omterm{def:m1:financing:repo}{décote} de 2~\% permettent $L = 50$~: c'est ainsi
que les fonds de relative value gagnent leur vie grâce à des écarts de prix de
quelques points de base (livre~9). Si les prêteurs portent la \omterm{def:m1:financing:repo}{décote} à 4~\%, un fonds
à la limite doit vendre \emph{la moitié} de ses actifs alors qu'aucun prix n'a bougé. Les \omterm{def:m1:financing:repo}{décotes}
sont fixées par les prêteurs, montent lorsque la volatilité augmente, et constituent le canal
par lequel un problème de financement devient un problème de marché.
\end{example}

\section{Prêt de titres et vente à découvert}

\begin{definition}[Vente à découvert]\label{def:m1:financing:short}
Une \emph{vente à découvert}\index{vente à découvert} est la vente d'un titre que le vendeur
ne possède pas. Pour livrer au règlement-livraison, le vendeur emprunte le titre, et
le rachète plus tard pour le restituer~: il gagne si le cours a baissé.
\end{definition}

\begin{definition}[Prêt de titres, taux de rebate, locate]\label{def:m1:financing:seclending}
Dans un \emph{prêt de titres}\index{prêt de titres} un propriétaire transfère un
titre à un emprunteur contre un collatéral, généralement du cash d'une valeur légèrement
supérieure à celle du titre, et reçoit en retour un titre équivalent sur demande. Le prêteur
place le cash et reverse une partie des intérêts à l'emprunteur au
\emph{taux de rebate}\index{taux de rebate}~: l'écart entre le taux d'intérêt de marché et le rebate
constitue les \emph{frais d'emprunt}\index{frais d'emprunt}, le véritable prix du prêt. Un
\emph{locate}\index{locate} est une confirmation du courtier, obtenue avant une \omterm{def:m1:financing:short}{vente à découvert},
selon laquelle le titre peut être emprunté à temps pour le règlement-livraison~; les règles
américaines en exigent un (\cref{ch:m1:stock-loan-and-short-selling}).
\end{definition}

\begin{omfigure}
\begin{tikzpicture}[font=\small,
  box/.style={draw=omDef,rounded corners=2pt,minimum width=2.7cm,minimum height=0.95cm,align=center,fill=omDef!4},
  flow/.style={-{Stealth[length=2mm]},thick}]
\node[box] (len) at (0,0) {Prêteur\\(un fonds de pension)};
\node[box] (pb) at (5.2,0) {Prime broker};
\node[box] (ss) at (10.4,0) {Vendeur à découvert};
\node[box] (mk) at (10.4,-2.6) {Marché};
\draw[flow,omDef] ([yshift=6pt]len.east) -- node[above,font=\footnotesize]{actions} ([yshift=6pt]pb.west);
\draw[flow,omThm] ([yshift=-6pt]pb.west) -- node[below,font=\footnotesize]{collatéral} ([yshift=-6pt]len.east);
\draw[flow,omDef] ([yshift=6pt]pb.east) -- node[above,font=\footnotesize]{actions} ([yshift=6pt]ss.west);
\draw[flow,omThm] ([yshift=-6pt]ss.west) -- node[below,font=\footnotesize]{produit} ([yshift=-6pt]pb.east);
\draw[flow,omDef] ([xshift=-6pt]ss.south) -- node[left,font=\footnotesize]{actions vendues} ([xshift=-6pt]mk.north);
\draw[flow,omThm] ([xshift=6pt]mk.north) -- node[right,font=\footnotesize]{produit de la vente} ([xshift=6pt]ss.south);
\end{tikzpicture}

\omcaption{Une \omterm{def:m1:financing:short}{vente à découvert}. Les actions (en bleu) passent d'un détenteur de long terme au
marché~; le cash (en rouge) revient et reste en garantie chez le prêteur, qui
le rémunère au \omterm{def:m1:financing:seclending}{taux de rebate}. Le vendeur à découvert ne détient jamais le
produit de la vente.}
\end{omfigure}

\begin{example}[Le carry d'une vente à découvert]\label{ex:m1:financing:carry}
Les taux d'intérêt sont de 4~\%. Une grande action très détenue se prête à des frais de 0.25~\%~:
le rebate vaut 3.75~\%, et le vendeur à découvert \emph{touche} 3.75~\% par an sur le produit
de la vente. Une action très demandée se prête à des frais de 15~\%~: le rebate vaut
$-11\%$, et le vendeur à découvert paie 11~\% par an pour le droit de rester vendeur. Dans
les deux cas, le vendeur à découvert verse aussi au prêteur tous les dividendes que l'action
distribue.
\end{example}

\section{Prime brokerage et marge}

\begin{definition}[Prime broker et réhypothécation]\label{def:m1:financing:pb}
Un \emph{prime broker}\index{prime broker} est la banque qui finance le portefeuille de titres d'un
\omterm{def:m1:what-a-trading-firm-does:hedge-fund}{hedge fund}~: elle prête du cash contre des longs, prête des titres pour les shorts, assure
le règlement-livraison et la conservation du portefeuille, et en rend compte.
La \emph{réhypothécation}\index{réhypothécation} désigne la réutilisation, par le prime broker, des
titres d'un client comme collatéral de ses propres emprunts, ce qui lui permet de financer les
prêts qu'il accorde.
\end{definition}

\begin{definition}[Appel de marge]\label{def:m1:financing:margincall}
Un \emph{appel de marge}\index{appel de marge} est la demande faite par un prêteur à un emprunteur
de rétablir l'excédent convenu du collatéral sur le prêt, en livrant du cash ou des titres, ou en
réduisant ses positions, en général dans la journée. Si l'appel n'est pas honoré, le prêteur
peut vendre le collatéral.
\end{definition}

\begin{dated}{2026-09}{Les règles de marge américaines en trois chiffres}\label{dat:m1:financing:rules}
La Regulation T de la Réserve fédérale fixe la marge \emph{initiale} sur un
achat d'actions cotées dans un compte sur marge à 50~\% (et à 150~\% pour une
\omterm{def:m1:financing:short}{vente à découvert}, dont le produit de la vente couvre 100~\%). La FINRA Rule 4210 fixe la
marge \emph{de maintien} à 25~\% de la valeur de marché pour les actions cotées en long. En
vertu de la SEC Rule 15c3-3, un courtier peut utiliser les titres d'un client à hauteur de 140~\%
du solde débiteur de ce client~; les titres au-delà de ce montant doivent rester en
possession ou sous le contrôle du courtier. Les exigences maison des \omterm{def:m1:financing:pb}{prime brokers}
et les régimes de marge de portefeuille s'ajoutent à ces règles.
\end{dated}

\begin{proposition}[Où se situe l'appel de marge]\label{prop:m1:financing:call}
Une position longue achetée à $P_0$ avec un ratio de \omterm{def:m1:clearing-and-settlement:margin}{marge initiale} $m_0$ (fonds propres
rapportés à la valeur) reçoit un \omterm{def:m1:financing:margincall}{appel de marge} de maintien au ratio de maintien $m$ lorsque le cours
passe sous
\[
P^\star = P_0\,\frac{1-m_0}{1-m}.
\]
\end{proposition}
\begin{proof}
Le prêt vaut $P_0(1-m_0)$ et ne varie pas avec le cours. Au cours $P$,
le ratio de fonds propres vaut $(P - P_0(1-m_0))/P$~; il suffit de l'égaler à $m$.
\end{proof}

\begin{example}[Cinquante et vingt-cinq]\label{ex:m1:financing:regt}
Avec $m_0 = 50\%$ et $m = 25\%$, $P^\star = \tfrac23 P_0$~: l'appel survient après
une baisse d'un tiers. Un client professionnel margé par son \omterm{def:m1:financing:pb}{prime broker}
à $m_0 = 15\%$ et $m = 10\%$ est appelé à $0.944\,P_0$, après une baisse de
5.6~\%.
\end{example}

\section{Financement synthétique}

\begin{definition}[Total return swap]\label{def:m1:financing:trs}
Dans un \emph{total return swap}\index{total return swap} sur actions, le dealer
verse au client le rendement total (variation du cours et dividendes) d'un titre sur
un montant notionnel, et le client verse un taux de financement majoré d'un spread sur le
même notionnel. Le client a l'économie d'une détention de l'action financée par
le dealer, qui détient les titres réels en guise de couverture~; le client dépose
une marge convenue dans le contrat.
\end{definition}

Le swap n'est qu'un financement sous un autre emballage, et l'emballage compte~: le
client n'apparaît sur aucun registre d'actionnaires~; la marge est celle que les deux
parties ont négociée, et non celle que prescrit la Regulation T~; et un client
qui travaille avec plusieurs dealers ne montre à chacun que sa part.

\begin{example}[Archegos]\label{ex:m1:financing:archegos}
En mars 2021, un family office, Archegos Capital Management, a fait défaut sur
les \omterm{def:m1:financing:margincall}{appels de marge} de ses dealers. Il avait constitué des positions concentrées sur
une poignée de titres au moyen de \omterm{def:m1:financing:trs}{total return swaps} conclus avec plusieurs banques, des positions,
que la plainte de la SEC a ensuite évaluées à 36~milliards de dollars. L'une de ces banques,
Credit Suisse, a perdu près de 5.5~milliards de dollars. Le rapport commandé par son conseil
d'administration a conclu que la banque avait accepté une marge sur swap de 7.5~\% ---
un \omterm{def:m1:financing:leverage}{effet de levier} supérieur à treize --- et que la marge était \emph{statique}, fixée sur le
cours auquel chaque swap avait été ouvert, de sorte que, à mesure que les titres montaient, la marge moyenne
détenue tombait à 6.9~\% de la valeur courante, et qu'un passage à une marge dynamique
n'avait pas été jugé prioritaire. Quand les titres ont chuté, chaque dealer
détenait les mêmes titres en guise de couverture, et chacun a dû les vendre dans le
flot des ventes des autres.
\end{example}

\section{Quand le financement se dérobe}

Combiner la \cref{prop:m1:financing:roe}(3) avec le fait que de grosses ventes font bouger
les cours (\cref{met:m1:the-sell-side:sqrt}) et le mécanisme de toute
crise à \omterm{def:m1:financing:leverage}{effet de levier} apparaît~: une baisse force à vendre, les ventes provoquent une nouvelle baisse,
qui force à vendre davantage.

\begin{omfigure}
\begin{tikzpicture}
\begin{axis}[width=10.5cm,height=4.8cm,
  xlabel={tour de ventes forcées},ylabel={baisse cumulée du cours (\%)},
  tick label style={font=\small},label style={font=\small},
  xmin=1,xmax=12,ymin=4,ymax=11,grid=major,grid style={gray!25},
  legend columns=3,legend style={font=\footnotesize,at={(0.5,-0.34)},anchor=north,draw=none,fill=none,/tikz/every even column/.append style={column sep=8pt}},
  table/col sep=comma]
\addplot[omMeth,very thick,mark=*,mark size=1.2pt] table[x=round,y=l3]{figdata/markets-1/06-financing/spiral.csv};
\addplot[omDef,very thick,mark=*,mark size=1.2pt] table[x=round,y=l5]{figdata/markets-1/06-financing/spiral.csv};
\addplot[omThm,very thick,mark=*,mark size=1.2pt] table[x=round,y=l8]{figdata/markets-1/06-financing/spiral.csv};
\legend{$L=3$,$L=5$,$L=8$}
\end{axis}
\end{tikzpicture}

\omcaption{Un choc de 5~\%, suivi de tours de ventes destinées à rétablir l'\omterm{def:m1:financing:leverage}{effet de levier}, chaque vente
faisant bouger le cours de 10~\% par unité d'actifs vendue (un impact illustratif).
Pour $L = 8$, le marché finit deux fois plus bas que le choc qui l'a déclenché.
Données~: le script du chapitre.}\label{fig:m1:financing:spiral}
\end{omfigure}

\begin{remark}[Trois accélérateurs]\label{rem:m1:financing:accelerants}
La spirale jouet de la \cref{fig:m1:financing:spiral} converge. Les vraies spirales sont
aggravées par trois éléments qu'elle omet. Les prêteurs relèvent les \omterm{def:m1:financing:repo}{décotes}
(la \cref{prop:m1:financing:haircut}) précisément quand les cours baissent. Les fonds qui détiennent
les mêmes positions sont touchés en même temps, si bien que le «~marché~» qui absorbe les ventes est
lui-même en train de vendre. Et les prêteurs qui redoutent la défaillance d'un emprunteur cessent de renouveler
les prêts au jour le jour, en bloc, transformant un \omterm{def:m1:financing:margincall}{appel de marge} en ruée.
Les pertes des fonds quantitatifs d'août 2007, la ruée sur le repo de 2008 et l'épisode de 2021 évoqué plus haut
sont les trois cas auxquels cette série revient.
\end{remark}

\section{Tutoriel~: le coût du carry d'un portefeuille long--short}\label{tut:m1:financing:carry}

\begin{tutorial}
\textbf{Objectif.} Calculer le coût annuel de financement d'un portefeuille d'actions long--short,
d'abord avec un compte de \omterm{def:m1:financing:pb}{prime broker} en espèces, puis avec des swaps.
\textbf{État final~:} deux nombres, 2.2~millions de dollars et 2.7~millions de dollars, et la
raison de l'écart.
\begin{tutsteps}
\item \textbf{Le portefeuille.} 300~millions de dollars en position longue, 200~millions en position courte, 100
millions de dollars de fonds propres. Conditions indicatives~: taux d'intérêt de 4~\%~; le \omterm{def:m1:financing:pb}{prime broker}
facture 0.50~\% de plus sur les soldes débiteurs et verse 0.30~\% de moins sur les produits
des \omterm{def:m1:financing:short}{ventes à découvert}~; les \omterm{def:m1:financing:seclending}{frais d'emprunt} moyens sont de 0.30~\%.
\item \textbf{Compte en espèces.} Le fonds emprunte $300 - 100 = 200$ millions de dollars.
\omcode{code/markets-1/06-financing/python/financing_demo.py}{60}{64}{Prime broker en espèces~: intérêts payés sur le débit, intérêts reçus sur les produits des ventes à découvert.}\label{lst:m1:financing:cash}
Coût~: $200 \times 4.5\% - 200 \times 3.4\% = 2.2$ millions de dollars par an.
\item \textbf{Swaps.} Le dealer finance la totalité du notionnel long et la
part du capital du fonds reste en espèces et rapporte le taux d'intérêt.
\omcode{code/markets-1/06-financing/python/financing_demo.py}{67}{72}{Le même portefeuille en swap.}\label{lst:m1:financing:swap}
Coût~: $300 \times 4.5\% - 200 \times 3.4\% - 100 \times 4\% = 2.7$ millions de dollars.
\item \textbf{Expliquer l'écart.} 0.5~million de dollars correspond au spread de 0.50~\% payé sur
des 100~millions de dollars de positions longues que le compte en espèces a financées avec l'argent
du fonds. Le financement synthétique facture le spread sur tout.
\end{tutsteps}
\textbf{Pour aller plus loin.} Placer un cinquième du portefeuille short en \omterm{def:m1:financing:seclending}{prêt de titres}
à des frais de 12~\%. Faire ensuite tomber le taux d'intérêt à zéro~: lequel des
deux coûts change-t-il, et pourquoi ni l'un ni l'autre ne change beaucoup~?
\end{tutorial}

\section{Construction~: le calculateur de financement}\label{bld:m1:financing:calc}

\begin{build}
\bfield{Rôle} Chaque nuit, la société de trading miniature calcule ce que lui coûtent
ses positions en carry et le porte au grand livre
du \cref{ch:m1:what-a-trading-firm-does} sous \texttt{financing}. \bfield{Interface} \texttt{FinancingTerms(debit\_spread, credit\_spread,
day\_count=360)}, une table par symbole des \omterm{def:m1:financing:seclending}{frais d'emprunt}, et
\texttt{accrue(date, positions, prices, equity, rate)} qui renvoie un montant signé
par ligne~: intérêts débiteurs, intérêts sur les produits des \omterm{def:m1:financing:short}{ventes à découvert}, \omterm{def:m1:financing:seclending}{frais d'emprunt}.
\bfield{Règles} Calculer sur les positions réglées, non sur les positions négociées. Les intérêts sont simples,
en jours exacts sur la base de comptage, de sorte qu'un vendredi couvre trois jours.
Les \omterm{def:m1:financing:seclending}{frais d'emprunt} se calculent sur la valeur de la position short évaluée à la clôture précédente.
Un \omterm{def:m1:financing:seclending}{frais d'emprunt} manquant pour une position short est une erreur, et non un zéro.
\bfield{Tests d'acceptation} \texttt{code/firm/financing/tests/}~: reproduit le
tutoriel, soit 2.2~millions de dollars sur une année de 360 jours~; un week-end
couvre trois jours~; un frais manquant lève une erreur.
\bfield{Pour aller plus loin} Ajouter des \omterm{def:m1:financing:repo}{décotes} par classe d'actifs et une requête
\texttt{max\_financing}~: combien de plus la société peut-elle acheter aujourd'hui~?
\end{build}

\begin{omsources}
\item Credit Suisse Group Special Committee of the Board of Directors,
\emph{Report on Archegos Capital Management}, 29 juillet 2021.
\item US Securities and Exchange Commission, \emph{SEC Charges Archegos and
its Founder with Massive Market Manipulation Scheme}, communiqué de presse 2022-70,
27 avril 2022.
\item 12 CFR 220.12 (Regulation T, supplement: margin requirements); FINRA
Rule 4210; SEA Rule 15c3-3 (FINRA interpretations handbook).
\item M.~Brunnermeier et L.~Pedersen, «~Market liquidity and funding
liquidity~», \emph{Review of Financial Studies} 22 (2009).
\item G.~Gorton et A.~Metrick, «~Securitized banking and the run on repo~»,
\emph{Journal of Financial Economics} 104 (2012).
\end{omsources}

\section{Exercices}

\begin{exercise}[$\star$]\label{exo:m1:financing:1}
Un fonds dispose de fonds propres de 250~millions de dollars et d'actifs de 1.5~milliard de dollars. Donner son
\omterm{def:m1:financing:leverage}{effet de levier}, la baisse de ses actifs qui l'anéantit, et son rendement des fonds propres
si les actifs rapportent $+6\%$ avec un financement à 4~\%.
\end{exercise}

\begin{exercise}[$\star$]\label{exo:m1:financing:2}
Une obligation valant 50~millions de dollars est financée en repo avec une \omterm{def:m1:financing:repo}{décote} de 3~\% et un \omterm{def:m1:financing:repo}{taux de repo}
de 4.2~\% pendant 7 jours (exact/360). Quel montant de liquidités est obtenu et quels
intérêts sont versés~?
\end{exercise}

\begin{exercise}[$\star$]\label{exo:m1:financing:3}
Les taux d'intérêt sont de 4~\%. Un vendeur à découvert est short de 20~millions de dollars d'un titre dont
les \omterm{def:m1:financing:seclending}{frais d'emprunt} sont de 6~\%. Que rapporte ou coûte la position par an en
financement, avant dividendes~?
\end{exercise}

\begin{exercise}[$\star\star$]\label{exo:m1:financing:4}
Une action est achetée à 80~dollars avec une \omterm{def:m1:clearing-and-settlement:margin}{marge initiale} de 40~\%~; le ratio de maintenance est
de 30~\%. À quel cours intervient l'\omterm{def:m1:financing:margincall}{appel de marge}~? L'action tombe à 60~dollars~: quelle
liquidité par action rétablit le ratio \emph{initial}~?
\end{exercise}

\begin{exercise}[$\star\star$]\label{exo:m1:financing:5}
Un fonds avec $L = 6$ et des actifs de 3~milliards de dollars perd 5~\% sur ses actifs.
Calculer son nouvel \omterm{def:m1:financing:leverage}{effet de levier}, et la vente nécessaire pour revenir à $L = 6$. Si ses
prêteurs n'acceptent plus que $L = 4$, combien doit-il vendre~?
\end{exercise}

\begin{exercise}[$\star\star$]\label{exo:m1:financing:6}
Un fonds de relative value fonctionne au levier maximal autorisé par une \omterm{def:m1:financing:repo}{décote}
de 2.5~\%. Les prêteurs relèvent la \omterm{def:m1:financing:repo}{décote} à 4~\%. Quelle fraction du portefeuille doit être
vendue~? La transaction rapporte 0.12~\% par an sur les actifs au-delà de son coût de financement~:
donner le rendement des fonds propres avant et après.
\end{exercise}

\begin{exercise}[$\star\star\star$]\label{exo:m1:financing:7}
\emph{Programmation.} Avec \texttt{spiral}, trouver par dichotomie l'impact par unité
vendue au-delà duquel un choc de 5~\% anéantit un fonds avec $L = 8$ en douze
tours. Comparer avec la valeur $1/(L-1)$ et expliquer le lien entre les deux.
\end{exercise}

\begin{exercise}[$\star\star\star$]\label{exo:m1:financing:8}
\emph{Trouver la faille.} Le rapport de risque d'un dealer indique au sujet d'un client en swaps~: «~La marge
détenue est de 7.5~\% du notionnel~; le plus fort mouvement quotidien du portefeuille au cours des
cinq dernières années a été de 6~\%~; l'exposition est entièrement couverte.~» La marge est
statique et les actions du client ont doublé depuis l'ouverture des swaps.
Identifier trois erreurs indépendantes.
\end{exercise}

\section{Problème~: un family office sur swaps}

\begin{problem}[{Problème du week-end --- levier à travers cinq dealers}]\label{pb:m1:financing:1}
Un family office dispose de 4~milliards de dollars de fonds propres. Par l'intermédiaire de \omterm{def:m1:financing:trs}{total return swaps} conclus avec
cinq dealers, à parts égales, il détient des positions longues sur huit actions dont la
valeur actuelle est de 28~milliards de dollars. Chaque dealer a prélevé une marge statique de 10~\% du
cours auquel ses swaps ont été ouverts~; les actions ont progressé de 60~\% en
moyenne depuis.

\medskip\noindent\textbf{Partie I --- Quel est son niveau de levier~?}
\begin{enumerate}
\item Donner le levier sur les valeurs actuelles et la baisse qui anéantit les
fonds propres.
\item Quel était le notionnel à l'ouverture des swaps, et quelle marge
les dealers détiennent-ils au total~?
\item Exprimer cette marge en pourcentage de la valeur actuelle.
\item Où se trouvent les 4~milliards de dollars restants du family office~? Pourquoi
est-il important pour les dealers de ne pas pouvoir les voir~?
\end{enumerate}

\medskip\noindent\textbf{Partie II --- Le carry.}
\begin{enumerate}[resume]
\item Les swaps coûtent le taux d'intérêt (4~\%) plus 0.45~\%. Donner le coût de financement annuel
en dollars.
\item La trésorerie du family office rapporte 4~\%. Donner son coût de carry net annuel et
l'exprimer en pourcentage des fonds propres.
\item Quel rendement annuel sur les actions lui faut-il pour atteindre le seuil de rentabilité~?
\end{enumerate}

\medskip\noindent\textbf{Partie III --- La chute.}
\begin{enumerate}[resume]
\item Les actions chutent de 10~\% en une journée. Donner la perte, les fonds propres restants et
le nouveau levier.
\item Chaque dealer appelle une \omterm{def:m1:clearing-and-settlement:margin}{marge de variation} égale à la perte sur sa tranche.
Quel montant de liquidités est appelé au total~? Le family office verse toute sa
trésorerie disponible, au prorata. Que reçoit chaque dealer, et ce qui lui reste dû~?
\item Le lendemain, les actions chutent de 8~\% supplémentaires. Reprendre la question 8.
\item Le family office est en défaut. Donner la valeur des actions de couverture d'un dealer
après les deux baisses, sa perte totale sur les swaps pour les deux jours, et
ce qu'il détient face à cette perte (marge plus liquidités reçues).
\item Un dealer vend ses actions dans les jours qui suivent à un cours moyen inférieur de 12~\%
à la clôture du second jour. Quelle est sa perte finale~?
\item Quelle aurait été sa perte s'il avait vendu immédiatement à 3~\% sous
cette clôture~?
\end{enumerate}

\medskip\noindent\textbf{Partie IV --- Ce qui aurait dû être différent.}
\begin{enumerate}[resume]
\item Avec une marge dynamique de 10~\% de la valeur \emph{courante}, quelle marge
les dealers auraient détenue avant la chute~?
\item Chaque position représente environ cinq jours de volume échangé sur l'action pour un
dealer seul. Avec une volatilité quotidienne de 3~\%, estimer la \omterm{def:m1:financing:repo}{décote} à 95~\% donnée par
la \cref{prop:m1:the-sell-side:discount} pour un dealer qui liquide seul, à un taux de 10~\%
de participation, puis comparer avec la marge détenue.
\item Pourquoi cette estimation reste-t-elle trop faible dans cette situation~?
\item Quelle information unique, si les dealers l'avaient partagée, aurait
modifié leur marge~?
\item Pourquoi le premier dealer à vendre a-t-il le moins perdu~?
\item Énoncer le \emph{résultat nommé}~: la baisse des actions qui efface
les fonds propres du family office, en pourcentage.
\item En une phrase~: quelle est la différence entre l'\omterm{def:m1:financing:leverage}{effet de levier} et la marge~?
\end{enumerate}
\end{problem}

\section{Questions d'entretien}

\begin{interviewq}[$\star$ \iqroles{trader, chercheur, banque}]\label{iq:m1:financing:1}
Décrivez-moi les flux de trésorerie d'une \omterm{def:m1:financing:short}{vente à découvert}, depuis le
\omterm{def:m1:financing:seclending}{locate} jusqu'au rachat.
\end{interviewq}

\begin{interviewq}[$\star$ \iqroles{banque, trader}]\label{iq:m1:financing:2}
Qu'est-ce qu'un repo, et pourquoi est-il jugé sûr pour le prêteur de liquidités~?
\end{interviewq}

\begin{interviewq}[$\star\star$ \iqroles{chercheur, trader}]\label{iq:m1:financing:3}
Un fonds a un \omterm{def:m1:financing:leverage}{effet de levier} de quatre et perd 5~\% sur ses actifs. Combien doit-il
vendre pour garder son \omterm{def:m1:financing:leverage}{effet de levier} constant~? Faites le calcul de tête.
\end{interviewq}

\begin{interviewq}[$\star\star$ \iqroles{chercheur}]\label{iq:m1:financing:4}
Votre backtest d'une stratégie long--short affiche 9~\% par an. Il ignore
le financement. Listez ce qu'il faut retrancher et donnez des ordres de grandeur.
\end{interviewq}

\begin{interviewq}[$\star\star$ \iqroles{banque}]\label{iq:m1:financing:5}
Pourquoi un client choisirait-il un \omterm{def:m1:financing:trs}{total return swap} plutôt que d'acheter les actions par l'intermédiaire de son
\omterm{def:m1:financing:pb}{prime broker}~? Pourquoi la banque préférerait-elle l'un ou l'autre~?
\end{interviewq}

\begin{interviewq}[$\star\star\star$ \iqroles{banque, chercheur}]\label{iq:m1:financing:6}
Vous êtes responsable des risques d'un \omterm{def:m1:financing:pb}{prime broker}. Le portefeuille d'un client est
concentré sur cinq actions, dans chacune desquelles il détient plusieurs jours de
volume, et vous soupçonnez qu'il détient les mêmes positions dans d'autres banques. Comment fixez-vous
sa marge~?
\end{interviewq}
